\documentclass[../main.tex]{subfiles}

\begin{document}

\textit{Chương này trình bày quá trình triển khai hệ thống: môi trường và cách cài đặt, cách các thành phần theo kiến trúc ở chương 3 được thi công, kết quả thi công các chức năng, các yêu cầu phi chức năng và bảo vệ hệ thống, dữ liệu thử nghiệm, kết quả kiểm thử, đánh giá và hướng dẫn sử dụng.}

%=====================================================================
\section{Môi trường và cài đặt hệ thống}\label{sec:c4-moi-truong}
%=====================================================================

Hệ thống được phát triển, kiểm thử và đo đạc trên máy cá nhân của sinh viên (AMD Ryzen 7 7735HS 16 luồng, 26~GB RAM, Ubuntu 26.04 LTS), sau đó triển khai thử nghiệm trên một máy chủ ảo Linux (6 vCPU Intel Xeon Platinum 8272CL 2,60~GHz, 16~GB RAM, Ubuntu 22.04.5 LTS). Máy chủ dùng chung với các dự án khác đã chiếm cổng 80/443, nên Apache được công bố ở cổng 8088 (HTTP) và 8443 (HTTPS). Hai môi trường chạy cùng một bộ cấu hình Docker Compose, chỉ khác giá trị biến môi trường. Yêu cầu tối thiểu để cài đặt là một máy Linux có Docker Engine, Docker Compose v2, tối thiểu 10~GB dung lượng trống và từ 4~GB RAM; Internet chỉ cần trong lúc cài đặt. Bảng~\ref{tab:c4-thu-vien} liệt kê phiên bản các công nghệ chính, được khóa cố định trong các tệp \path{uv.lock}, \path{package-lock.json} để mọi lần cài đặt cho cùng một kết quả.

{\renewcommand{\arraystretch}{1.15}
\begin{longtable}{|p{2.3cm}|p{6.2cm}|p{6.1cm}|}
\caption{Công nghệ và thư viện chính}\label{tab:c4-thu-vien}\\
\hline
\textbf{Thành phần} & \textbf{Công nghệ (phiên bản)} & \textbf{Vai trò} \\
\hline
\endfirsthead
\hline
\textbf{Thành phần} & \textbf{Công nghệ (phiên bản)} & \textbf{Vai trò} \\
\hline
\endhead
Web API & Python 3.12, FastAPI 0.141, Uvicorn 0.54, Pydantic 2.13, SQLAlchemy 2.0.54, Alembic 1.20, PyJWT 2.15, bcrypt 5.0, python-docx 1.2, PyMuPDF 1.28 & Khung API, ORM và migration, ký JWT, băm mật khẩu, trích xuất văn bản DOCX/PDF \\
\hline
Dịch vụ AI & sentence-transformers 6.1, PyTorch 2.14 (CPU), underthesea 9.5, faiss-cpu 1.15, NumPy 2.5; mô hình \path{bkai-foundation-models/vietnamese-bi-encoder} & Mã hóa câu thành vector~\cite{bkaibiencoder}, tách từ tiếng Việt~\cite{underthesea}, chỉ mục vector~\cite{faiss2021}, chỉ mục vân tay văn bản (winnowing~\cite{schleimer2003}) tự cài đặt \\
\hline
Client Web & React 18.3, TypeScript 5.9, Vite 7.3, Ant Design 5.29, TanStack Query 5.104, React Router 7.18, Zustand 5 & Giao diện, quản lý dữ liệu từ máy chủ, định tuyến, trạng thái phiên \\
\hline
CSDL & MySQL 8.4 & Hai schema \texttt{thesis\_core}, \texttt{thesis\_ai}, bảng mã \texttt{utf8mb4\_vietnamese\_ci} \\
\hline
Triển khai & Apache HTTP Server 2.4, Docker Compose & Reverse proxy, HTTPS, phục vụ client tĩnh; chạy các container \\
\hline
Kiểm thử & pytest 8, Chrome headless, Bash, curl & Kiểm thử đơn vị, tích hợp, giao diện, end-to-end \\
\hline
\end{longtable}
}

Mỗi thành phần được đóng gói thành một image Docker, toàn bộ hệ thống được mô tả trong tệp \path{deploy/docker-compose.yml}~\cite{dockercompose} (Hình~\ref{fig:c4-docker-compose}). Sơ đồ này hiện thực hai vùng của sơ đồ triển khai mức thiết kế (Hình~\ref{fig:c3-trien-khai}): vùng công khai chỉ có container \textit{apache}, công bố cổng 80 và 443, phục vụ bản build tĩnh của client và chuyển tiếp \texttt{/api/}; vùng nội bộ là mạng Docker \texttt{ntma-internal}, gồm \textit{api} (Web API, cổng 8001), \textit{ai} (dịch vụ AI, cổng 8002, hệ thống tệp chỉ đọc, chạy ngoại tuyến), \textit{mysql} và hai service chạy một lần \textit{ai-db-init}, \textit{ai-migrate} tạo và nâng cấp schema của dịch vụ AI. Web API và dịch vụ AI không công bố cổng ra máy chủ; MySQL chỉ mở cổng trên \texttt{127.0.0.1} cho quản trị viên. Dữ liệu cần giữ qua các lần cập nhật nằm trong ba volume: \texttt{mysql-data}, \texttt{uploads-data} (nằm ngoài \texttt{DocumentRoot} của Apache) và \texttt{ai-models}.

\begin{figure}[H]
    \centering
    \includegraphics[width=0.95\textwidth]{chuong_4/so-do-docker-compose.png}
    \caption{Sơ đồ cài đặt hệ thống bằng Docker Compose}
    \label{fig:c4-docker-compose}
\end{figure}

Một máy chủ MySQL chứa hai schema tách biệt, mỗi schema có tài khoản CSDL riêng: \texttt{thesis\_core} là dữ liệu nghiệp vụ của Web API, được tạo và nâng cấp bằng 19 migration Alembic~\cite{alembic}; \texttt{thesis\_ai} là dữ liệu của dịch vụ AI (văn bản nguồn đã chuẩn hóa, các câu, vector embedding), có ba migration riêng. Tài khoản của dịch vụ AI bị từ chối nếu truy cập \texttt{thesis\_core}. Mọi cột văn bản dùng \texttt{utf8mb4\_vietnamese\_ci}~\cite{mysqlutf8mb4}, nên tìm kiếm và sắp xếp tiếng Việt đúng thứ tự từ điển.

Việc cài mới hoặc nâng cấp được gói trong một lệnh \texttt{bash deploy/install.sh}, chạy lại an toàn: tập lệnh kiểm tra điều kiện, sinh tệp biến môi trường với các bí mật ngẫu nhiên và chứng chỉ TLS tự ký, tải trước trọng số mô hình (khoảng 518~MiB, bước duy nhất cần Internet), khởi động và chờ mọi service ở trạng thái healthy, kiểm tra migration và tài khoản, rồi nạp 300 đề tài kho và 25 đồ án nguồn vào dịch vụ AI. Hai lần cài từ đầu với volume hoàn toàn mới mất 216~s (máy cũ, có sẵn bộ đệm build Docker) và 545~s (máy hiện tại, không có bộ đệm), đều kết thúc thành công không cần thao tác tay.

%=====================================================================
\section{Thi công các thành phần theo kiến trúc}
%=====================================================================

\textbf{Web API} là ứng dụng FastAPI~\cite{fastapi} sở hữu toàn bộ nghiệp vụ, gồm 17 nhóm router với tiền tố \texttt{/api}, mỗi nhóm ứng với một module nghiệp vụ; tài liệu OpenAPI được sinh tự động tại \path{/api/docs}. Mã nguồn tổ chức theo tầng: \path{routers/} khai báo endpoint và vai trò được phép, \path{services/} chứa nghiệp vụ, \path{models/} là các lớp thực thể SQLAlchemy, \path{schemas/} là các model Pydantic. Các chuyển trạng thái quan trọng (nộp, duyệt đề cương) dùng câu lệnh \texttt{UPDATE} có điều kiện trên trạng thái cũ hoặc \texttt{SELECT ... FOR UPDATE}, nên khi hai người thao tác đồng thời chỉ một thao tác thành công, thao tác còn lại nhận mã 409. Tệp tải lên được trích xuất văn bản bằng python-docx (DOCX) hoặc PyMuPDF (PDF). Với PDF, Web API dựng lại từng dòng từ vị trí của từng ký tự, vì PyMuPDF 1.28 không chèn dấu cách sau các chữ có dấu thuộc khối Unicode U+1E00--U+1EFF khi PDF (thường sinh từ LaTeX) không chứa ký tự dấu cách thật, làm các âm tiết bị dính liền (``hệthống''). Trên 104 PDF đồ án thật, số lần xuất hiện ở dạng dính liền của 6 cụm từ thông dụng giảm từ 39\,663 xuống 0; đổi lại, trích xuất chậm hơn khoảng 3,4 lần (trung vị 0,35 giây mỗi tệp).

\textbf{Client Web} là ứng dụng React một trang~\cite{react} viết bằng TypeScript~\cite{typescript}, build bằng Vite~\cite{vite} thành tệp tĩnh do Apache phục vụ, dùng bộ thành phần Ant Design~\cite{antd}. Mỗi route khai báo danh sách vai trò được phép và menu được sinh từ cùng danh sách, nên mỗi vai trò chỉ thấy chức năng của mình; quyết định cho phép luôn nằm ở Web API. Mọi lời gọi API đi qua TanStack Query~\cite{tanstackquery}; access token chỉ giữ trong bộ nhớ, refresh token nằm trong cookie \texttt{HttpOnly}.

\textbf{Dịch vụ AI} là một ứng dụng FastAPI riêng, sở hữu schema \texttt{thesis\_ai}, cung cấp năm endpoint \texttt{/ingest}, \texttt{/check-title}, \texttt{/check-document}, \texttt{/reindex} và \texttt{/health}; mọi endpoint trừ \texttt{/health} yêu cầu header \texttt{X-API-Key}. Khi khởi động, dịch vụ nạp mô hình \texttt{vietnamese-bi-encoder}~\cite{bkaibiencoder} từ volume và dựng hai chỉ mục FAISS trong bộ nhớ, một cho tên đề tài, một cho các câu của tài liệu. Chỉ mục dùng tích vô hướng trên vector đã chuẩn hóa (tương đương cosine) và cho phép thêm, xóa theo định danh, nên nạp một nguồn mới chỉ mã hóa đúng nguồn đó. Việc mã hóa được chia thành lô con 32 câu, nên một lời gọi kiểm tra tên đến giữa lúc đang kiểm tra một báo cáo dài chỉ phải chờ tối đa một lô con.

\textbf{Giao tiếp giữa các thành phần.} Trình duyệt chỉ nói chuyện với Apache qua HTTPS; Apache chuyển các yêu cầu \texttt{/api/} sang Web API qua \texttt{mod\_proxy}~\cite{apachemodproxy}. Web API là thành phần duy nhất gọi dịch vụ AI, bằng REST JSON trong mạng nội bộ, chỉ gửi văn bản đã trích xuất. Kiểm tra trùng lặp tài liệu chạy bất đồng bộ: Web API lưu tệp, tạo bản ghi kết quả ở trạng thái \textit{Đang kiểm tra}, trả mã 202 ngay rồi gọi dịch vụ AI trong \texttt{BackgroundTasks}; client thăm dò cho tới khi nhận \textit{Hoàn tất} hoặc \textit{Lỗi}. Nếu dịch vụ AI tạm ngừng, việc nộp vẫn thành công và kết quả có nút \textit{Kiểm tra lại}. Sau mỗi lần ghi đề tài, bản ghi kho hoặc tài liệu, Web API đẩy dữ liệu sang \texttt{/ingest} để kho so sánh luôn cập nhật.

%=====================================================================
\section{Kết quả thi công các chức năng}
%=====================================================================

Bảng~\ref{tab:c4-tong-hop-uc} tổng hợp các chức năng đã thi công theo từng UC; các hình sau đó là ảnh chụp giao diện thật, chụp tự động bằng Chrome headless ở khung 1440$\times$900 điểm ảnh trên một bản cài đặt riêng với dữ liệu hư cấu (mục~\ref{sec:c4-du-lieu}), không chứa dữ liệu của người thật.

{\renewcommand{\arraystretch}{1.15}
\begin{longtable}{|p{2.6cm}|p{9.4cm}|p{2.6cm}|}
\caption{Tổng hợp kết quả thi công theo UC}\label{tab:c4-tong-hop-uc}\\
\hline
\textbf{UC} & \textbf{Chức năng đã thi công} & \textbf{Vai trò} \\
\hline
\endfirsthead
\hline
\textbf{UC} & \textbf{Chức năng đã thi công} & \textbf{Vai trò} \\
\hline
\endhead
UC0 -- Chức năng chung & Đăng nhập, đăng xuất, làm mới phiên; quản lý tài khoản, đợt đồ án, phân công GVHD, danh mục định hướng; cấu hình ngưỡng; thông báo; kiểm tra trùng tên đề tài, trùng lặp tài liệu và kiểm tra đạo văn một tệp tải lên (mục~\ref{sec:c4-phi-chuc-nang}); trang giới thiệu, tổng quan theo vai trò, trang 403 & Mọi vai trò \\
\hline
UC1 -- Quản lý đề tài & GVHD tạo và giao đề tài cho sinh viên trong lớp hướng dẫn, kèm cảnh báo trùng tên theo ngữ nghĩa; xem, lọc danh sách đề tài; sinh viên xem đề tài được giao và lịch sử & GV, SV, Bộ môn \\
\hline
UC2 -- Quản lý đề cương & Sinh viên nộp đề cương (DOCX/PDF), kiểm tra tệp, tên tệp, định hướng và trùng lặp tự động; GVHD duyệt hoặc yêu cầu sửa, nộp lên Bộ môn, đổi tên tệp theo mẫu; Bộ môn phê duyệt bằng bản PDF đã ký hoặc trả lại, xuất ZIP; lịch sử phiên bản & SV, GV, Bộ môn \\
\hline
UC3 -- Quản lý tiến độ & Bộ môn thiết lập mốc; sinh viên nộp báo cáo theo mốc; GVHD nhận xét, đánh giá; Bộ môn theo dõi toàn khoa; lịch sử nộp của sinh viên & SV, GV, Bộ môn \\
\hline
UC4 -- Quản lý bảo vệ & Nộp báo cáo cuối kèm kiểm tra trùng lặp; lập hội đồng, xếp lịch; thành viên hội đồng chấm điểm; Bộ môn chốt kết quả; sinh viên xem lịch và kết quả & SV, GV, Bộ môn \\
\hline
UC5 -- Quản lý kho đồ án & Tra cứu, xem chi tiết, tải báo cáo; tự lưu kho khi chốt bảo vệ; thêm, sửa, xóa bản ghi; import CSV các khóa cũ; trạng thái chỉ mục AI & Mọi vai trò xem; Bộ môn ghi \\
\hline
\end{longtable}
}

\textbf{UC0. Chức năng chung.} Sau khi đăng nhập, mỗi vai trò có một trang tổng quan riêng. Trang của giảng viên tập trung vào việc cần xử lý (sinh viên chưa có đề tài, đề cương chờ duyệt, báo cáo tiến độ chờ nhận xét), mỗi việc có nút đi thẳng tới thao tác tương ứng (Hình~\ref{fig:c4-gv-tong-quan}). Bộ môn quản lý tài khoản, đợt đồ án, phân công GVHD, danh mục định hướng và các ngưỡng cảnh báo trùng lặp (Hình~\ref{fig:c4-nguong}); mỗi ngưỡng hiển thị giá trị đề xuất từ bước hiệu chỉnh (mục~\ref{sec:c4-kiem-thu}) và chỉ áp dụng cho các lần kiểm tra sau. Hình được chụp trước khi bổ sung ngưỡng tổng nên mới có ba ngưỡng. Người dùng mở chức năng không thuộc vai trò của mình sẽ thấy trang 403.

\begin{figure}[H]
    \centering
    \includegraphics[width=0.8\textwidth]{chuong_4/c0-gv-tong-quan.png}
    \caption{Tổng quan của giảng viên: việc cần xử lý}
    \label{fig:c4-gv-tong-quan}
\end{figure}

\begin{figure}[H]
    \centering
    \includegraphics[width=0.8\textwidth]{chuong_4/c0-bm-nguong.png}
    \caption{Cấu hình ngưỡng cảnh báo trùng lặp}
    \label{fig:c4-nguong}
\end{figure}

\textbf{UC1. Quản lý đề tài.} Sinh viên thống nhất đề tài với GVHD ngoài hệ thống, sau đó GVHD giao đề tài trên hệ thống cho sinh viên trong lớp hướng dẫn. Khi GVHD nhập tên đề tài, client gọi API kiểm tra trùng tên; Hình~\ref{fig:c4-canh-bao-ten} là trường hợp tên mới diễn đạt lại một đề tài trong kho, với tỷ lệ tương đồng 81,92\%, trên ngưỡng 67,8\%. Cảnh báo chỉ mang tính gợi ý: GVHD có thể sửa tên hoặc chọn \textit{Vẫn giao đề tài}. Danh sách đề tài (Hình~\ref{fig:c4-ds-de-tai}) lọc được theo đợt, trạng thái, GVHD, định hướng; mỗi vai trò chỉ thấy đề tài trong phạm vi của mình.

\begin{figure}[H]
    \centering
    \includegraphics[width=0.8\textwidth]{chuong_4/c1-gv-canh-bao-trung-ten.png}
    \caption{Cảnh báo trùng tên đề tài khi GVHD giao đề tài}
    \label{fig:c4-canh-bao-ten}
\end{figure}

\begin{figure}[H]
    \centering
    \includegraphics[width=0.8\textwidth]{chuong_4/c1-bm-danh-sach-de-tai.png}
    \caption{Danh sách đề tài của đợt}
    \label{fig:c4-ds-de-tai}
\end{figure}

\textbf{UC2. Quản lý đề cương.} Đề cương đi qua quy trình duyệt hai cấp: sinh viên nộp, GVHD duyệt và nộp lên Bộ môn, Lãnh đạo bộ môn phê duyệt bằng bản PDF đề cương đã ký. Sau khi sinh viên nộp, phiên bản mới có trạng thái \textit{Đang kiểm tra} rồi chuyển sang \textit{Hoàn tất} kèm tỷ lệ trùng (Hình~\ref{fig:c4-sv-nop-de-cuong}).

\begin{figure}[H]
    \centering
    \includegraphics[width=0.75\textwidth]{chuong_4/c2-sv-nop-de-cuong.png}
    \caption{Sinh viên nộp đề cương và xem kết quả kiểm tra}
    \label{fig:c4-sv-nop-de-cuong}
\end{figure}

Màn hình duyệt của GVHD đặt thông tin đề cương, danh sách kiểm tra trước khi nộp lên Bộ môn, nội dung tệp và kết quả kiểm tra trùng lặp trên cùng một trang (Hình~\ref{fig:c4-gv-vuot-nguong}: đề cương trùng 48\% với một đồ án trong kho, vượt ngưỡng 10,4\%; tên tệp không theo mẫu và đề tài chưa gắn định hướng). Nút \textit{Xem chi tiết} mở trang kết quả chi tiết (Hình~\ref{fig:c4-chi-tiet-trung-lap}): bên trái là danh sách đồ án trùng kèm tỷ lệ; bên phải là từng đoạn trùng đặt cạnh đoạn nguồn, cụm từ trùng được tô sáng, mỗi đoạn có nhãn mức trùng và loại trùng (nguyên văn, sửa nhẹ hoặc ngữ nghĩa).

\begin{figure}[H]
    \centering
    \includegraphics[width=0.92\textwidth,height=0.45\textheight,keepaspectratio]{chuong_4/c2-gv-duyet-vuot-nguong.png}
    \caption{GVHD xem đề cương có tỷ lệ trùng vượt ngưỡng}
    \label{fig:c4-gv-vuot-nguong}
\end{figure}

Trước khi nộp lên Bộ môn, hệ thống kiểm tra tên tệp đúng quy tắc \textit{[Mã định hướng].[MSSV].[Họ tên].[Lớp].pdf} và mã định hướng khớp định hướng của đề tài (điều kiện bắt buộc), còn kết quả trùng lặp chỉ để Bộ môn tham khảo; nếu lỗi duy nhất là tên tệp, GVHD có thể bấm \textit{Đổi tên theo mẫu}. Ở cấp Bộ môn, Lãnh đạo bộ môn xem lại các mục kiểm tra, tải đề cương về, ký ngoài hệ thống rồi tải bản PDF đã ký lên để phê duyệt, hoặc trả lại kèm nhận xét bắt buộc (Hình~\ref{fig:c4-bm-cho-duyet}). Hệ thống chỉ yêu cầu tệp tải lên là PDF, không đọc hay xác minh chữ ký trong tệp. Bản đã ký được lưu vào thư mục đề cương của đợt, và Bộ môn tải được toàn bộ thư mục dưới dạng một tệp ZIP để nộp về khoa.

\begin{figure}[H]
    \centering
    \includegraphics[width=0.88\textwidth]{chuong_4/c2-chi-tiet-trung-lap.png}
    \caption{Chi tiết kết quả kiểm tra trùng lặp: đoạn nộp cạnh đoạn nguồn}
    \label{fig:c4-chi-tiet-trung-lap}
\end{figure}

\begin{figure}[H]
    \centering
    \includegraphics[width=0.92\textwidth,height=0.45\textheight,keepaspectratio]{chuong_4/c2-bm-cho-phe-duyet.png}
    \caption{Bộ môn phê duyệt đề cương bằng bản PDF đã ký}
    \label{fig:c4-bm-cho-duyet}
\end{figure}

\textbf{UC3. Quản lý tiến độ thực hiện.} Bộ môn thiết lập các mốc tiến độ của đợt. Sinh viên nộp báo cáo theo từng mốc; hệ thống ghi nhận nộp đúng hạn hay trễ hạn và hiển thị đánh giá của GVHD (Hình~\ref{fig:c4-sv-tien-do}). GVHD gửi nhận xét kèm đánh giá \textit{Đạt}, \textit{Cần cải thiện} hoặc \textit{Chưa đạt}; Bộ môn theo dõi tiến độ toàn khoa và lọc được sinh viên trễ mốc.

\begin{figure}[H]
    \centering
    \includegraphics[width=0.8\textwidth]{chuong_4/c3-sv-bao-cao-tien-do.png}
    \caption{Sinh viên nộp báo cáo tiến độ theo mốc}
    \label{fig:c4-sv-tien-do}
\end{figure}

\textbf{UC4. Quản lý bảo vệ.} Sinh viên có đề cương đã được phê duyệt nộp báo cáo cuối, được kiểm tra trùng lặp như đề cương. Bộ môn lập hội đồng (3 đến 5 thành viên, đúng một chủ tịch), chọn ngày, giờ, phòng và thêm sinh viên vào lịch (Hình~\ref{fig:c4-hoi-dong}). Mỗi thành viên hội đồng chấm điểm và nhận xét từng sinh viên, kèm tóm tắt kiểm tra trùng lặp của báo cáo cuối (Hình~\ref{fig:c4-cham-diem}). Khi đủ điểm, Bộ môn chốt kết quả; điểm bảo vệ là trung bình cộng điểm của các thành viên, sau khi chốt không sửa được và đồ án tự động được lưu vào kho.

\begin{figure}[H]
    \centering
    \includegraphics[width=0.8\textwidth]{chuong_4/c4-bm-hoi-dong.png}
    \caption{Lập hội đồng và lịch bảo vệ}
    \label{fig:c4-hoi-dong}
\end{figure}

\begin{figure}[H]
    \centering
    \includegraphics[width=0.8\textwidth]{chuong_4/c4-gv-cham-diem.png}
    \caption{Thành viên hội đồng chấm điểm}
    \label{fig:c4-cham-diem}
\end{figure}

\textbf{UC5. Quản lý kho dữ liệu đồ án.} Kho lưu các đồ án đã bảo vệ, gồm đồ án tự lưu khi chốt bảo vệ và đồ án các khóa cũ do Bộ môn import từ tệp CSV. Mọi vai trò tra cứu được theo khóa, GVHD và từ khóa; cột trạng thái chỉ mục AI cho biết bản ghi đã vào kho so sánh hay chưa (Hình~\ref{fig:c4-kho-tim-kiem}). Mỗi dòng import thành công được đẩy sang dịch vụ AI, nên tên đề tài cũ tham gia ngay vào các lần kiểm tra trùng tên sau đó.

\begin{figure}[H]
    \centering
    \includegraphics[width=0.8\textwidth]{chuong_4/c5-kho-tim-kiem.png}
    \caption{Tra cứu kho đồ án}
    \label{fig:c4-kho-tim-kiem}
\end{figure}

%=====================================================================
\section{Thi công các yêu cầu phi chức năng}\label{sec:c4-phi-chuc-nang}
%=====================================================================

Ba yêu cầu phi chức năng gắn trực tiếp với mục tiêu của đề tài là kiểm tra trùng tên đề tài, kiểm tra trùng lặp tài liệu bằng AI và phân quyền theo vai trò. Bảo mật được trình bày ở mục~\ref{sec:c4-bao-ve}, kết quả đo ở mục~\ref{sec:c4-kiem-thu}.

\subsection{Kiểm tra trùng lặp bằng AI}\label{sec:c4-ai-trung-lap}

\textbf{Trùng tên đề tài.} Kiểm tra chạy đồng bộ ngay khi GVHD nhập xong tên đề tài. Web API chuẩn hóa tên về dạng Unicode NFC và gửi tên cùng ngưỡng tên hiện hành sang \texttt{/check-title} của dịch vụ AI, thời gian chờ 4 giây. Dịch vụ AI mã hóa tên thành vector 768 chiều, tìm 15 vector gần nhất trong chỉ mục tên đề tài (chứa cả đề tài của đợt hiện hành và tên đồ án trong kho), gộp theo nguồn, lọc theo ngưỡng và trả tối đa 5 nguồn có cosine cao nhất. Kết quả chỉ là cảnh báo: khi dịch vụ AI không phản hồi, GVHD vẫn giao được đề tài. Kết quả tại thời điểm giao (kể cả ngưỡng đã dùng) được lưu kèm đề tài, và tên vừa giao được đẩy ngay vào chỉ mục.

\textbf{Trùng lặp tài liệu.} Đề cương và báo cáo cuối được kiểm tra để chỉ ra ba thông tin: tỷ lệ trùng của tài liệu, các đồ án bị trùng kèm tỷ lệ với từng đồ án, và các đoạn trùng cụ thể. Dịch vụ AI xử lý văn bản đã trích xuất như sau:
\begin{enumerate}
    \item \textit{Loại phần mẫu chung.} Bìa, lời cam đoan, mục lục, danh mục hình và bảng, tài liệu tham khảo, phụ lục, tiêu đề chương mục, chú thích, số trang bị loại khỏi phép so sánh vì mọi đồ án đều giống nhau. Trên đồ án thật, sau khi bổ sung quy tắc cho dải chấm mục lục của PDF sinh từ LaTeX và dòng giá trị của bìa, số câu còn chứa dải chấm giảm từ 7\,176 xuống 0, số mảnh bìa còn sót giảm từ 106 xuống 7. Để quy tắc không thành chỗ giấu nội dung, các mục mẫu chỉ được bỏ ở phần đầu và phần cuối tài liệu, có trần độ dài, và kết quả luôn ghi số dòng đã loại.
    \item \textit{Tách câu và mã hóa.} Văn bản được chuẩn hóa NFC (kể cả chữ ư mà font TCVN3 mã hóa thành ký tự khác), nối các dòng bị PDF ngắt rồi tách thành câu; câu ngắn hơn 6 âm tiết bị bỏ qua, câu dài hơn 120 âm tiết được cắt nhỏ, nhờ vậy cùng một nội dung ở dạng DOCX và PDF cho cùng các câu.
    \item \textit{Hai tín hiệu.} Tín hiệu ngữ nghĩa là cosine với câu gần nghĩa nhất trong chỉ mục tài liệu. Tín hiệu nguyên văn dùng vân tay văn bản~\cite{schleimer2003,sorokina2006}: mỗi câu được băm thành các cụm 6 âm tiết liên tiếp, kho chỉ giữ các vân tay chọn bằng winnowing (cửa sổ 4 cụm) trong một chỉ mục ngược, bảo đảm mọi đoạn chung dài từ 9 âm tiết trở lên đều được tìm thấy. Điểm nguyên văn là \textit{độ phủ}: phần âm tiết của câu nằm trong một đoạn chung dài từ 9 âm tiết với một câu của cùng một đồ án nguồn, lấy hợp trên mọi câu ứng viên của nguồn đó. Chép cả câu cho 100\%, sửa một từ giữa câu còn khoảng 97\%, ghép hai câu của cùng một nguồn vẫn được khoảng 100\%. Cụm có mặt trong từ 4 đề tài khác nhau trở lên bị coi là phổ biến và không được tính; các phiên bản trước của chính đề tài đó không tham gia so sánh.
    \item \textit{Gắn cờ.} Câu bị gắn cờ \textit{nguyên văn} khi độ phủ đạt 56\%; nếu không, câu bị gắn cờ \textit{ngữ nghĩa} khi cosine đạt ngưỡng đoạn trùng (67,7\%) và câu thuộc một chuỗi ít nhất 2 câu liên tiếp cùng trùng một nguồn, để loại các câu đơn lẻ trùng ý ngẫu nhiên.
    \item \textit{Tổng hợp.} Tỷ lệ trùng tổng là số ký tự của các câu bị gắn cờ chia cho tổng số ký tự đã so sánh; tỷ lệ trùng với từng đồ án tính tương tự. Các câu bị gắn cờ liên tiếp cùng nguồn được ghép thành đoạn trùng (tối đa 6 câu), mang nhãn \textit{trùng ngữ nghĩa} khi dưới một nửa số ký tự bị gắn cờ nguyên văn; nếu không, \textit{trùng sửa nhẹ} khi tỷ lệ âm tiết của đoạn khớp liên tục với chữ của đoạn nguồn dưới 0,99, ngược lại \textit{trùng nguyên văn}. Nhãn chỉ để giải thích cho người xem, không làm đổi tỷ lệ.
\end{enumerate}

Web API đánh dấu \textit{vượt ngưỡng} khi tỷ lệ trùng với \textit{một} đồ án nguồn cao nhất đạt ngưỡng cảnh báo tài liệu (bắt việc chép phần lớn một đồ án), hoặc khi tỷ lệ trùng tổng đạt ngưỡng tổng (bắt bài ghép từ nhiều đồ án, mỗi đồ án một ít). Giá trị hiệu chỉnh của ngưỡng tổng (14,8\%) cao hơn ngưỡng thứ nhất (10,4\%) vì câu tự viết cùng lĩnh vực có thể khớp rải rác với nhiều đồ án; kết quả ghi rõ lý do vượt ngưỡng. Trên máy chủ đang vận hành, từ ngày 29/10/2026 Bộ môn đặt cả hai ngưỡng là 25\%; vì mỗi câu bị gắn cờ chỉ thuộc về một nguồn nên tỷ lệ theo một đồ án không vượt tỷ lệ tổng, tức hệ thống cảnh báo khi tỷ lệ trùng tổng đạt 25\%. Tài liệu chỉ được đẩy vào chỉ mục sau khi kết quả \textit{Hoàn tất}, nên không bao giờ tự trùng với chính nó. Mọi lỗi trong tác vụ nền, hoặc bản ghi quá 300 giây chưa xong, đều chuyển sang \textit{Lỗi}; văn bản dài hơn 800\,000 ký tự (khoảng 230 trang) không được gửi sang dịch vụ AI.

\textbf{Kiểm tra đạo văn một tệp tải lên.} Theo yêu cầu của giảng viên hướng dẫn, quy trình trên được mở thành một chức năng độc lập: giảng viên và Bộ môn tải một tệp PDF hoặc DOCX bất kỳ lên màn hình \textit{Kiểm tra đạo văn} để so với kho đồ án và các tài liệu đã nộp. Tệp được kiểm tra định dạng và trích xuất như khi nộp, rồi đi qua đúng tác vụ nền của tài liệu nộp; kết quả hiển thị trên cùng màn hình chi tiết với tỷ lệ tổng, đồ án trùng nhiều nhất và từng đoạn trùng cạnh đoạn nguồn. Tệp không được lưu và không được đưa vào chỉ mục, nên không ảnh hưởng các lần kiểm tra bài của sinh viên sau đó; kết quả chỉ người tải lên xem được, và mỗi người có tối đa hai lần kiểm tra đang chạy. Tệp không trích xuất được chữ, như PDF dạng ảnh quét, bị từ chối ngay.

\subsection{Phân quyền}

Hệ thống có ba vai trò: Sinh viên, Giảng viên (gồm GVHD và thành viên hội đồng) và Bộ môn. Lớp quyết định là Web API, kiểm tra ở ba mức. Thứ nhất, mọi route dưới \texttt{/api} khai báo vai trò được phép qua dependency \texttt{require\_role}, gọi sai vai trò nhận mã 403; chỉ tám route được để công khai (đăng nhập, làm mới phiên, đăng xuất, kiểm tra tình trạng, tài liệu OpenAPI), và một kiểm thử tự động báo lỗi nếu có route mới quên khai báo vai trò. Thứ hai, mỗi yêu cầu đọc lại tài khoản từ CSDL: tài khoản bị khóa, vai trò không khớp hoặc token phát ra trước lần đăng xuất gần nhất đều nhận mã 401, nên khóa tài khoản và đăng xuất có hiệu lực ngay. Thứ ba, các hàm nghiệp vụ kiểm tra quyền trên từng đối tượng theo phạm vi dữ liệu (Bảng~\ref{tab:c4-phan-quyen}).

{\renewcommand{\arraystretch}{1.15}
\begin{longtable}{|p{4.6cm}|p{2.9cm}|p{3.4cm}|p{3.4cm}|}
\caption{Quyền chính của từng vai trò}\label{tab:c4-phan-quyen}\\
\hline
\textbf{Nhóm chức năng} & \textbf{Sinh viên} & \textbf{Giảng viên} & \textbf{Bộ môn} \\
\hline
\endfirsthead
\hline
\textbf{Nhóm chức năng} & \textbf{Sinh viên} & \textbf{Giảng viên} & \textbf{Bộ môn} \\
\hline
\endhead
Tài khoản, đợt, phân công GVHD, định hướng, ngưỡng & Xem GVHD của mình & Xem lớp hướng dẫn & Toàn quyền \\
\hline
Đề tài & Xem đề tài của mình & Tạo, giao cho SV mình hướng dẫn & Xem toàn khoa \\
\hline
Đề cương & Nộp, nộp lại & Duyệt, yêu cầu sửa, nộp lên Bộ môn, đổi tên tệp & Phê duyệt bằng bản PDF đã ký, trả lại, xuất ZIP \\
\hline
Tiến độ & Nộp báo cáo theo mốc & Nhận xét SV mình hướng dẫn & Thiết lập mốc, theo dõi toàn khoa \\
\hline
Bảo vệ & Nộp báo cáo cuối, xem lịch, điểm & Chấm điểm trong hội đồng của mình & Lập hội đồng, xếp lịch, chốt điểm \\
\hline
Kho đồ án & Tra cứu, xem & Tra cứu, xem & Thêm, sửa, xóa, import, đánh chỉ mục lại \\
\hline
Kết quả kiểm tra trùng lặp & Của mình & Của SV mình hướng dẫn hoặc trong hội đồng & Toàn khoa \\
\hline
\end{longtable}
}

%=====================================================================
\section{Bảo vệ hệ thống}\label{sec:c4-bao-ve}
%=====================================================================

Các biện pháp bảo vệ chính gồm:
\begin{itemize}
    \item \textit{Apache và HTTPS:} Apache là điểm vào duy nhất; cổng 80 chuyển hướng mọi yêu cầu sang HTTPS, cổng 443 chỉ chấp nhận TLS 1.2 và 1.3 với chứng chỉ tự ký RSA 2048 bit sinh lúc cài đặt, thay được bằng chứng chỉ do tổ chức cấp chứng chỉ phát hành. Apache ẩn phiên bản, tắt \texttt{TRACE}, không liệt kê thư mục, giới hạn thân yêu cầu 25~MB, gửi các header bảo mật và chính sách nội dung (CSP), và chỉ có một đường proxy \texttt{/api/} sang Web API.
    \item \textit{Xác thực:} mật khẩu băm bằng bcrypt với hệ số chi phí 12, có phép băm giả khi tài khoản không tồn tại để thời gian phản hồi không tiết lộ tài khoản có thật; sau 5 lần sai trong 15 phút từ cùng địa chỉ IP với cùng tên đăng nhập, các lần thử tiếp nhận mã 429. Access token hạn 30 phút chỉ giữ trong bộ nhớ của trang; refresh token hạn 8 giờ nằm trong cookie \texttt{HttpOnly}, \texttt{Secure}, \texttt{SameSite=Strict}; đăng xuất làm mọi token cũ mất hiệu lực.
    \item \textit{Tệp tải lên:} kiểm tra theo nội dung chứ không theo phần mở rộng (PDF bắt đầu bằng \texttt{\%PDF-}, DOCX là tệp ZIP hợp lệ có \path{word/document.xml}), tối đa 20~MB, chống tệp nén ``bom'' (giải nén không quá 25~MB); tệp được lưu dưới tên do máy chủ sinh, ngoài \texttt{DocumentRoot}, chỉ tải được qua API có kiểm tra quyền. Bản đề cương đã ký cũng qua đúng các bước này, cộng điều kiện phải là PDF; hệ thống không xác minh chữ ký, hệ quả được nêu ở phần kết luận chương.
    \item \textit{Dịch vụ AI và CSDL:} chỉ nằm trong mạng Docker nội bộ; dịch vụ AI yêu cầu khóa nội bộ, chạy với hệ thống tệp chỉ đọc; mỗi dịch vụ có tài khoản CSDL riêng chỉ có quyền trên schema của mình. Các bí mật được sinh ngẫu nhiên lúc cài đặt, chỉ nằm trong tệp biến môi trường có quyền đọc chỉ cho chủ sở hữu; danh sách sinh viên, giảng viên thật không nằm trong kho mã nguồn.
\end{itemize}

%=====================================================================
\section{Dữ liệu thử nghiệm}\label{sec:c4-du-lieu}
%=====================================================================

Dữ liệu thử nghiệm được chuẩn bị theo các chỉ tiêu NFR-17 đến NFR-19 ở chương 2 và được nạp hoàn toàn bằng tập lệnh, chạy lại không sinh bản ghi trùng: 580 tài khoản sinh viên và 70 giảng viên (danh sách thật dùng cho triển khai), 300 đề tài kho khóa K18--K22 thuộc 8 cụm chủ đề, 80 truy vấn tên đề tài có nhãn (40 diễn đạt lại một đề tài trong kho, 40 cùng lĩnh vực khác đối tượng) và 50 tệp tài liệu gồm 25 cặp A/B (7 chép nguyên văn, 7 diễn đạt lại, 6 cùng lĩnh vực khác đề tài, 5 không liên quan; 17 cặp dùng hiệu chỉnh, 8 cặp giữ riêng để kiểm chứng). Kho đề tài, truy vấn và tài liệu là dữ liệu tổng hợp do bộ sinh \path{tools/datagen} tạo với hạt giống cố định, không chứa văn bản của đồ án thật, mô phỏng khuôn mẫu chung của đồ án CNTT tiếng Việt để thử được cả bước loại phần mẫu chung.

Kho đồ án trên máy chủ triển khai đã nạp 106 trong 128 đồ án thật của khoa; 22 đồ án còn lại phần lớn chỉ có tệp nguồn LaTeX. Dữ liệu thật chưa được gán nhãn trùng hay không trùng nên ngưỡng vẫn được hiệu chỉnh trên dữ liệu tổng hợp; để đo trên văn bản thật, mục~\ref{sec:c4-pan} dựng thêm một bộ đánh giá kiểu PAN từ báo cáo đồ án thật. Bộ dữ liệu này chỉ lưu cục bộ, không đưa vào kho mã nguồn.

%=====================================================================
\section{Kiểm thử và đánh giá}\label{sec:c4-kiem-thu}
%=====================================================================

Hệ thống được kiểm thử ở ba mức: kiểm thử tự động toàn hệ thống, thử nghiệm riêng cho mô hình AI, và đo các chỉ tiêu phi chức năng ở chương 2. Mọi số liệu được sinh tự động bởi tập lệnh và lưu kèm nhật ký trong \path{docs/chuong-4-trien-khai/}.

\subsection{Kiểm thử tự động}

Tập lệnh \texttt{deploy/verify-all.sh -{}-report} chạy tuần tự 14 bước trên stack đang chạy (qua Apache HTTPS, với dịch vụ AI và mô hình thật); các bước ghi dữ liệu dùng đợt và tài khoản kiểm thử riêng và tự dọn sau khi chạy. Bảng~\ref{tab:c4-verify-all} so sánh lần chạy 14/10/2026, trước khi có luồng duyệt đề cương hai cấp, với lần 21/10/2026 trên mã nguồn hiện tại, chạy trên một stack riêng với danh sách người dùng hư cấu.

{\renewcommand{\arraystretch}{1.15}
\begin{longtable}{|p{0.6cm}|p{5.3cm}|p{2.1cm}|p{2.1cm}|p{3.6cm}|}
\caption{Kết quả kiểm thử tự động toàn hệ thống}\label{tab:c4-verify-all}\\
\hline
\textbf{\#} & \textbf{Bước} & \textbf{14/10/2026} & \textbf{21/10/2026} & \textbf{Ghi chú lần 21/10} \\
\hline
\endfirsthead
\hline
\textbf{\#} & \textbf{Bước} & \textbf{14/10/2026} & \textbf{21/10/2026} & \textbf{Ghi chú lần 21/10} \\
\hline
\endhead
1 & Tình trạng container, mô hình ngoại tuyến & Đạt & Không đạt & Do môi trường đo \\
\hline
2 & Dữ liệu thử nghiệm đủ số lượng & Đạt & Đạt & \\
\hline
3 & Smoke test qua HTTPS & Đạt (26) & Đạt (26) & \\
\hline
4 & Kiểm thử Web API (pytest) & 54 ca đạt & 67 ca đạt & Thêm ca luồng Bộ môn phê duyệt \\
\hline
5 & Kiểm thử dịch vụ AI (pytest) & 8 ca đạt & 8 ca đạt & \\
\hline
6 & E2E UC1: phân công GVHD, giao đề tài & Đạt (15) & Đạt (15) & \\
\hline
7 & E2E UC2--UC5 & Đạt (67) & Không đạt (11) & Kịch bản chưa cập nhật \\
\hline
8 & E2E dịch vụ AI thật & Đạt (31) & Đạt (31) & \\
\hline
9 & Hiệu năng, bảo mật dịch vụ AI & Đạt & Đạt & \\
\hline
10 & Yêu cầu hệ thống qua HTTPS & Đạt (14) & Đạt (14) & \\
\hline
11 & Đối chiếu ERD với schema & Đạt & Không đạt & ERD chưa cập nhật \\
\hline
12 & Không còn dữ liệu kiểm thử & Đạt & Đạt & \\
\hline
13 & Hiển thị màn hình chính ở 1440 px & Đạt (25) & Không đạt & Kịch bản chưa cập nhật \\
\hline
14 & Trạng thái demo không đổi & Đạt & Đạt & \\
\hline
\end{longtable}
}

\noindent Số trong ngoặc là số kiểm tra đã qua. Lần 14/10 đạt cả 14 bước trong 336 giây. Lần 21/10 đạt 10/14 bước trong 430 giây, và cả bốn bước không đạt đều do tập lệnh kiểm thử hoặc tài liệu chưa theo kịp thay đổi, không phát hiện lỗi của chức năng: bước 1 gọi cố định cổng 443 trong khi stack riêng dùng cổng 18443; bước 7 vẫn chờ đề cương chuyển thẳng sang \textit{Đã duyệt} sau khi GVHD duyệt (luồng hai cấp được kiểm thử ở mức API bằng các ca pytest mới của bước 4); bước 11 dùng tệp ERD cũ (với ERD đã cập nhật cho chương 3, schema khớp hoàn toàn); bước 13 dùng kịch bản chụp màn hình cũ (kịch bản mới mở được 35/35 màn hình ở 1440 px). Sau khi thay tín hiệu nguyên văn bằng vân tay văn bản và thêm ngưỡng tổng (22/10/2026), 68/68 ca của Web API và 15/15 ca của dịch vụ AI đều đạt; sau các thay đổi ngày 24/11/2026 (đánh giá kiểu PAN, kiểm tra tệp tải lên), 75/75 ca của Web API và 32/32 ca của dịch vụ AI đều đạt. Toàn bộ 14 bước chưa được chạy lại sau các thay đổi này.

\subsection{Thử nghiệm mô hình AI}

\textbf{Chọn mô hình embedding.} Hai mô hình \texttt{vietnamese-bi-encoder} và \texttt{multilingual-e5-base} được so sánh trên 80 cặp văn bản gán nhãn tay (30 cặp tên đề tài, 30 cặp câu, 20 cặp đoạn), theo quy tắc viết ra trước khi chạy: chỉ chọn \texttt{multilingual-e5-base} khi AUC của nó vượt ít nhất 0,05 trên \textit{cả} tập câu và tập đoạn. Trên tập câu hai mô hình bằng nhau (AUC 1,0000; tập đoạn 0,9479 so với 1,0000; tên đề tài 0,8472 so với 0,8750), nên \texttt{vietnamese-bi-encoder} được chọn. Mô hình này còn nhỏ hơn (khoảng 515~MiB so với 1,06~GiB), mã hóa 300 câu nhanh hơn (4,84~s so với 6,30~s) và trải điểm rộng hơn hẳn (chênh lệch cosine trung bình giữa cặp trùng và không trùng 0,24--0,47, so với 0,06--0,09), giúp ngưỡng ít nhạy với sai số hiệu chỉnh.

\textbf{Hiệu chỉnh ngưỡng.} Các ngưỡng được chọn trên tập hiệu chỉnh (17 cặp tài liệu, 56 truy vấn tên đề tài), rồi kiểm chứng trên tập giữ riêng (8 cặp, 24 truy vấn). Ngưỡng tên, đoạn và nguyên văn được chọn theo chỉ số Youden J trên đường ROC; ngưỡng tài liệu và ngưỡng tổng lấy trung điểm vì hai lớp tách biệt hoàn toàn (Bảng~\ref{tab:c4-nguong}). Độ dài cụm của vân tay được chọn cùng lúc theo AUC của tín hiệu nguyên văn với cụm 3, 4, 5 và 6 âm tiết (lần lượt 0,9859; 0,9962; 0,9969; 0,9974), nên hệ thống dùng cụm 6 âm tiết.

{\renewcommand{\arraystretch}{1.15}
\begin{longtable}{|p{4.3cm}|p{1.7cm}|p{3.1cm}|p{1.6cm}|p{3cm}|}
\caption{Ngưỡng sau hiệu chỉnh}\label{tab:c4-nguong}\\
\hline
\textbf{Ngưỡng} & \textbf{Giá trị} & \textbf{Phương pháp} & \textbf{AUC} & \textbf{TPR / FPR} \\
\hline
\endfirsthead
\hline
\textbf{Ngưỡng} & \textbf{Giá trị} & \textbf{Phương pháp} & \textbf{AUC} & \textbf{TPR / FPR} \\
\hline
\endhead
Tên đề tài (cosine) & 67,8\% & Youden J & 0,9758 & 0,893 / 0,036 \\
\hline
Đoạn trùng ngữ nghĩa (cosine từng câu) & 67,7\% & Youden J & 0,7983 & 0,829 / 0,351 \\
\hline
Nguyên văn (độ phủ vân tay từng câu) & 56,0\% & Youden J & 0,9974 & 0,978 / 0,019 \\
\hline
Tài liệu (tỷ lệ theo một đồ án) & 10,4\% & Trung điểm & 1,0000 & 1,000 / 0,000 \\
\hline
Tài liệu (tỷ lệ tổng mọi đồ án) & 14,8\% & Trung điểm & 1,0000 & 1,000 / 0,000 \\
\hline
\end{longtable}
}

\textbf{Kiểm chứng.} Bảng~\ref{tab:c4-kiem-chung} cho độ chính xác ở ba mức. Ở mức tài liệu, mọi cặp chép nguyên văn, diễn đạt lại và không liên quan đều được xử lý đúng trên cả hai tập; cặp duy nhất sai là P19 (cùng lĩnh vực, tập kiểm chứng), trùng ngữ nghĩa 22,81\% với một đồ án khác cùng chủ đề y tế trong kho nên bị cảnh báo. Ở mức câu, câu chép nguyên văn được phát hiện đầy đủ, nhưng câu diễn đạt lại chỉ được phát hiện 40--56\%; tài liệu diễn đạt lại vẫn bị cảnh báo đúng vì đủ nhiều câu trùng tập trung vào một nguồn. Ba cổng chất lượng đặt trước đều đạt: mọi cặp chép nguyên văn bị cảnh báo, mọi cặp không liên quan dưới ngưỡng, độ chính xác cặp trên tập kiểm chứng 0,875 (yêu cầu $\geq 0{,}85$).

{\renewcommand{\arraystretch}{1.15}
\begin{longtable}{|p{6.8cm}|p{3.8cm}|p{3.8cm}|}
\caption{Độ chính xác kiểm tra trùng lặp trên tập hiệu chỉnh và tập kiểm chứng}\label{tab:c4-kiem-chung}\\
\hline
\textbf{Chỉ số} & \textbf{Tập hiệu chỉnh} & \textbf{Tập kiểm chứng} \\
\hline
\endfirsthead
\hline
\textbf{Chỉ số} & \textbf{Tập hiệu chỉnh} & \textbf{Tập kiểm chứng} \\
\hline
\endhead
\multicolumn{3}{|l|}{\textit{Mức tài liệu (cặp A/B, cảnh báo đúng kỳ vọng)}} \\
\hline
Chép nguyên văn & 5/5 & 2/2 \\
\hline
Diễn đạt lại & 5/5 & 2/2 \\
\hline
Cùng lĩnh vực, khác đề tài & 4/4 & 1/2 \\
\hline
Không liên quan & 3/3 & 2/2 \\
\hline
\textbf{Tất cả} & \textbf{17/17 (1,000)} & \textbf{7/8 (0,875)} \\
\hline
\multicolumn{3}{|l|}{\textit{Mức câu (tỷ lệ câu bị gắn cờ)}} \\
\hline
Câu chép nguyên văn & 183/183 (1,000) & 110/110 (1,000) \\
\hline
Câu diễn đạt lại & 125/222 (0,563) & 22/55 (0,400) \\
\hline
Câu tự viết bị gắn cờ nhầm & 60/1064 (0,056) & 66/616 (0,107) \\
\hline
\multicolumn{3}{|l|}{\textit{Mức tên đề tài (truy vấn có nhãn)}} \\
\hline
Độ chính xác chung & 52/56 (0,929) & 22/24 (0,917) \\
\hline
Diễn đạt lại được cảnh báo & 25/28 (0,893) & 11/12 (0,917) \\
\hline
Cùng lĩnh vực không bị cảnh báo & 27/28 (0,964) & 11/12 (0,917) \\
\hline
\end{longtable}
}

\textbf{Bài ghép từ nhiều nguồn.} Bộ dữ liệu không có loại bài này, nên một kịch bản được dựng thêm, không dùng để hiệu chỉnh: vào bài tự viết của 11 cặp cùng lĩnh vực hoặc không liên quan, chèn bốn khối câu chép nguyên văn từ bốn đồ án khác trong kho, chiếm 19,9--22,5\% bài mới. Cả 284 câu chèn đều được gắn cờ nguyên văn, nhưng ở 9 bài tỷ lệ trùng với một đồ án cao nhất chỉ 5,2--6,1\%, nên quy tắc chỉ xét từng nguồn chỉ cảnh báo 2/11 bài. Khi thêm ngưỡng tổng, cả 11/11 bài bị cảnh báo, với tỷ lệ tổng 20,3--39,3\%.

\textbf{Ngưỡng đang áp dụng trên máy chủ.} Tính lại trên cùng số liệu với ngưỡng 25\% của Bộ môn: mọi cặp chép nguyên văn (tỷ lệ tổng từ 44\% trở lên) vẫn bị cảnh báo và cặp P19 (23,96\%) không còn bị cảnh báo nhầm, nhưng ba cặp diễn đạt lại có tỷ lệ tổng khoảng 21\% không còn bị cảnh báo. Độ chính xác mức tài liệu thành 16/17 trên tập hiệu chỉnh và 6/8 trên tập kiểm chứng; kịch bản ghép bốn nguồn cảnh báo 6/11 bài. Ngưỡng cao hơn giảm số cảnh báo cần giảng viên xem lại, đổi lại bỏ sót một phần bài diễn đạt lại hoặc ghép với tỷ lệ trùng dưới 25\%.

\textbf{Vai trò của hai bước lọc.} Khi tắt bước loại phần mẫu chung, tỷ lệ trùng tổng của 11 bài cùng lĩnh vực hoặc không liên quan tăng trung bình 8,02 điểm phần trăm và 3 bài vượt ngưỡng tổng. Giữ bước này nhưng tắt bộ lọc cụm phổ biến thì tỷ lệ tổng lên 7,6--35,0\% (P17 từ 8,41\% lên 25,66\%) và 6/11 bài vượt ngưỡng tổng, vì các câu theo khuôn chung của lĩnh vực bị tính là chép nguyên văn. Hai bước lọc vì vậy bổ sung cho nhau, và là điều kiện để dùng được ngưỡng tổng.

\subsection{Đánh giá kiểu PAN trên đồ án thật}\label{sec:c4-pan}

Để đo hệ thống trên văn bản thật và biết nó bắt được kiểu chép nào, bỏ sót kiểu nào, một bộ đánh giá được dựng theo khung PAN-PC~\cite{potthast2010}: chèn các đoạn chép đã biến đổi vào báo cáo đồ án thật, ghi vị trí ký tự của từng đoạn, rồi đo phần ký tự hệ thống phát hiện được. Bản đề xuất thuật toán của giảng viên hướng dẫn cũng gợi ý cách này, vì chưa có bộ dữ liệu chuẩn công khai cho tiếng Việt.

\textbf{Dữ liệu và cách chấm.} Bộ v1 dùng 105 báo cáo đồ án thật khóa K19--K21, trích xuất bằng đúng hàm của Web API: 75 đồ án làm kho nguồn, 30 làm tài liệu kiểm tra, trong đó 20 được chèn đạo văn và 10 để nguyên. Mỗi tài liệu được chèn nhận 7 đoạn (3--8 câu thân bài liền nhau, 60--250 âm tiết), mỗi đoạn lấy từ một nguồn khác nhau và mang một kiểu biến đổi khác nhau, tổng cộng 140 ca, 20 ca mỗi kiểu. Bảy kiểu được đặt tên theo bảng phân loại ReParaphrased~\cite{zhou2024} và kiểu che giấu ngẫu nhiên của PAN~\cite{potthast2010}: \path{nguyen_van} (chép nguyên văn), \path{sua_nhe} (thay khoảng 15--25\% số cụm bằng từ đồng nghĩa, bỏ vài hư từ, đổi dấu câu), \path{dao_cau} (đảo thứ tự câu), \path{gop_cau} (gộp từng cặp câu bằng từ nối), \path{tach_cau} (tách câu dài hơn 25 âm tiết tại dấu phẩy), \path{xao_tron} (xáo từ trong khoảng 30\% cửa sổ 3--4 từ, xóa và chèn khoảng 5\% số từ) và \path{llm} (một mô hình ngôn ngữ lớn (Claude) viết lại theo yêu cầu kiểu sinh viên: diễn đạt khác, giữ nguyên ý và thuật ngữ, tránh chuỗi từ 8 âm tiết trở lên trùng bản gốc). Bộ v2 được dựng lại từ cùng các tệp PDF với cùng hạt giống, bằng hàm trích xuất PDF mới và bộ lọc phần mẫu chung mới: còn 104 đồ án (74 kho nguồn, 30 kiểm tra), các đoạn chèn được chọn lại và các đoạn \path{llm} được viết lại mới. Hệ thống chạy đúng pipeline của dịch vụ AI với ngưỡng đoạn 67,7\% và ngưỡng nguyên văn 56,0\%; một phát hiện là một chuỗi tối đa các câu liên tiếp bị gắn cờ cùng một nguồn. Theo PAN-PC~\cite{potthast2010}, \textit{recall} là tỷ lệ ký tự của mỗi ca được phát hiện phủ, \textit{precision} là tỷ lệ ký tự của mỗi phát hiện nằm trong một ca, \textit{granularity} là số phát hiện trung bình trên mỗi ca được phát hiện, và \textit{plagdet} là F1 chia cho $\log_2(1 + \text{granularity})$.

{\renewcommand{\arraystretch}{1.15}
\begin{longtable}{|>{\raggedright\arraybackslash}p{4.4cm}|p{3cm}|p{3cm}|>{\raggedright\arraybackslash}p{3.6cm}|}
\caption{Kết quả đánh giá kiểu PAN trên bộ v1 (sau bản sửa độ phủ) và bộ v2}\label{tab:c4-pan-ket-qua}\\
\hline
\textbf{Chỉ số} & \textbf{Bộ v1} & \textbf{Bộ v2} & \textbf{Bộ v1: recall do tín hiệu nguyên văn} \\
\hline
\endfirsthead
\hline
\textbf{Chỉ số} & \textbf{Bộ v1} & \textbf{Bộ v2} & \textbf{Bộ v1: recall do tín hiệu nguyên văn} \\
\hline
\endhead
\multicolumn{4}{|l|}{\textit{Tổng thể (140 ca)}} \\
\hline
Precision & 0,2294 & 0,2912 & \\
\hline
Recall & 0,8544 & 0,8907 & \\
\hline
Granularity & 1,000 & 1,037 & \\
\hline
Plagdet & 0,3617 & 0,4275 & \\
\hline
Số phát hiện / không trùng ca nào & 536 / 410 & 464 / 325 & \\
\hline
Ca phát hiện được & 126/140 & 134/140 & \\
\hline
\multicolumn{4}{|l|}{\textit{Theo kiểu biến đổi: recall (số ca phát hiện được)}} \\
\hline
\path{nguyen_van} & 0,950 (19/20) & 0,982 (20/20) & 0,950 \\
\hline
\path{sua_nhe} & 1,000 (20/20) & 0,993 (20/20) & 0,941 \\
\hline
\path{dao_cau} & 1,000 (20/20) & 0,995 (20/20) & 1,000 \\
\hline
\path{gop_cau} & 0,928 (20/20) & 0,924 (20/20) & 0,928 \\
\hline
\path{tach_cau} & 0,950 (19/20) & 0,959 (20/20) & 0,948 \\
\hline
\path{xao_tron} & 0,590 (14/20) & 0,589 (16/20) & 0,098 \\
\hline
\path{llm} & 0,563 (14/20) & \textbf{0,793 (18/20)} & 0,000 \\
\hline
\end{longtable}
}

\textbf{Kết quả theo kiểu biến đổi.} Lần chạy đầu trên bộ v1 cho precision 0,2298, recall 0,8503, plagdet 0,3619, và cho thấy một lỗ hổng: câu ghép từ hai câu nguồn chỉ được khoảng 50\% độ phủ với mỗi câu nguồn, dưới ngưỡng nguyên văn. Sau khi độ phủ được tính theo từng nguồn (mục~\ref{sec:c4-ai-trung-lap}), phần recall do tín hiệu nguyên văn của kiểu gộp câu tăng từ 0,590 lên 0,928 (Bảng~\ref{tab:c4-pan-ket-qua}). Chép nguyên văn, sửa nhẹ, đảo câu, gộp câu và tách câu được bắt gần đủ, gần như toàn bộ nhờ vân tay; trên bộ v2 mọi ca của năm kiểu này đều được phát hiện. Hai ca ``bỏ sót'' của bộ v1 lấy từ một đồ án nguồn chứa cùng một đoạn hai lần, hệ thống khớp với lần xuất hiện thứ nhất còn ca được ghi ở lần thứ hai, nên đó là đặc điểm của dữ liệu. Hai kiểu khó là xáo trộn ngẫu nhiên và viết lại bằng mô hình ngôn ngữ lớn. Đoạn \path{llm} chỉ được bắt bằng tín hiệu ngữ nghĩa, phù hợp nhận xét rằng embedding huấn luyện sẵn dùng nguyên rồi so cosine còn yếu với câu diễn đạt lại, và cần huấn luyện có giám sát trên các cặp câu có nhãn để cải thiện~\cite{iqbal2025}; trên bộ v2, kiểu này đạt recall 0,793, chỉ 2/20 ca không được phát hiện.

\textbf{Mức tài liệu.} Không tài liệu sạch nào bị cảnh báo nhầm, ở cả hai bộ và cả hai cặp ngưỡng (tỷ lệ trùng tổng 1,3--7,3\% ở bộ v1, 0,03--6,50\% ở bộ v2). Tài liệu được chèn có tỷ lệ tổng 3,3--17,1\% (bộ v1) và 3,6--16,75\% (bộ v2, trừ một tài liệu ngoại lai). Với ngưỡng hiệu chỉnh 10,4\% / 14,8\%, 2/20 (bộ v1) và 3/20 (bộ v2) tài liệu được chèn bị cảnh báo; với ngưỡng 25\% / 25\%, không tài liệu nào ở bộ v1 và chỉ tài liệu ngoại lai ở bộ v2. Điều này không có nghĩa là hệ thống bỏ sót các bài chép: phần chèn chỉ gồm 7 đoạn trong cả một đồ án (khoảng 5--10\% nội dung), lại chia cho 7 nguồn. Ngưỡng theo tỷ lệ đo \textit{chép bao nhiêu}, không đo \textit{có chép hay không}; mọi đoạn phát hiện vẫn hiện trong danh sách đoạn trùng để giảng viên xem. Tài liệu ngoại lai có tỷ lệ tổng 51,04\% (26,77\% với một nguồn), phần lớn không thuộc các đoạn được chèn: đó là một đồ án thật khóa K21 có nhiều đoạn trùng nguyên văn dài tới khoảng 2\,000 ký tự với hai đồ án thật cùng nhóm đề tài trong kho nguồn. Hệ thống đã làm lộ ra việc dùng lại văn bản thật giữa các đồ án; còn đó là khuôn mẫu chung của nhóm hay là sao chép thì cần người xem lại để kết luận.

\textbf{Giới hạn của phép đo.} Precision thấp (0,23 và 0,29) vì nền là đồ án thật: các phát hiện không chồng lên ca được chèn nào gồm nội dung thật sự giống nhau giữa các đồ án cùng chủ đề, kiến thức giáo khoa chung và phần mẫu chung lọt qua bộ lọc cũ; trong 10 phát hiện dài nhất của bộ v2, 8 là trùng nguyên văn và 9 thuộc tài liệu ngoại lai nói trên, nên precision đo được chỉ là cận dưới của precision thật. Hai bộ \textit{không} tạo thành một phép so sánh có đối chứng: ngoài bản sửa trích xuất và bộ lọc mẫu chung, các đoạn chèn đã được chọn lại, kho nguồn ít hơn một đồ án, và các đoạn \path{llm} mới bám độ dài bản gốc hơn (0,91--1,08 lần, so với 1,03--1,23 lần ở bộ v1); bản sửa trích xuất là giải thích hợp lý nhất cho mức tăng của các kiểu dựa vào tín hiệu ngữ nghĩa, nhưng phần đóng góp này chưa được đo tách riêng. Cuối cùng, nhãn \textit{trùng sửa nhẹ} ban đầu dựa trên điểm nguyên văn (ngưỡng 96,3\%), nhưng kiểm thử đầu-cuối cho thấy độ phủ vân tay hụt ở mép câu làm đoạn chép y hệt bị gắn nhầm; nhãn được đổi sang tỷ lệ khớp chữ với đoạn nguồn (ngưỡng 0,99), trên 40 ca của bộ v2 tách hoàn toàn hai kiểu (chép nguyên văn đều 1,00, sửa nhẹ 0,875--0,976). Nhãn chỉ để hiển thị, không đổi cảnh báo.

\subsection{Đánh giá theo yêu cầu phi chức năng}

Bảng~\ref{tab:c4-nfr} đối chiếu kết quả lần chạy ngày 21/10/2026 với 19 yêu cầu phi chức năng ở chương 2. Các phép đo thời gian chạy trên máy phát triển với dữ liệu thử nghiệm đầy đủ (308 vector tên đề tài, 2\,991 vector câu tài liệu), khi tín hiệu nguyên văn còn dùng MinHash quét toàn kho; tín hiệu vân tay thay thế nhanh hơn nhiều: trên kho tổng hợp 150\,000 câu, so khớp nguyên văn cho một báo cáo 1\,500 câu mất 0,09 giây, so với 34,5 giây khi quét bằng MinHash. Trên máy chủ triển khai chỉ chạy các phép đo không ghi dữ liệu, mỗi phép đo gọi 20 lần qua Apache HTTPS.

{\renewcommand{\arraystretch}{1.15}
\begin{longtable}{|p{1.5cm}|p{4.4cm}|p{6.6cm}|p{1.7cm}|}
\caption{Kết quả đánh giá yêu cầu phi chức năng (21/10/2026)}\label{tab:c4-nfr}\\
\hline
\textbf{Mã} & \textbf{Chỉ tiêu} & \textbf{Kết quả đo} & \textbf{Đánh giá} \\
\hline
\endfirsthead
\hline
\textbf{Mã} & \textbf{Chỉ tiêu} & \textbf{Kết quả đo} & \textbf{Đánh giá} \\
\hline
\endhead
NFR-01 & Mật khẩu băm bcrypt, chi phí $\geq 12$, không ghi log & 651/651 tài khoản có tiền tố \texttt{\$2b\$12\$}; mật khẩu khởi tạo xuất hiện 0 lần trong log & Đạt \\
\hline
NFR-02 & Access token 30 phút, refresh token 8 giờ, đăng xuất thu hồi token & Token hết hạn 401, còn hạn 200; sau đăng xuất, làm mới phiên trả 401 & Đạt \\
\hline
NFR-03 & 5 lần sai trong 15 phút thì chặn, mã 429 & Lần đăng nhập sai thứ 6 nhận 429 & Đạt \\
\hline
NFR-04 & Kiểm tra tệp theo nội dung, tối đa 20 MB & Tệp đổi phần mở rộng, tệp quá cỡ bị từ chối & Đạt \\
\hline
NFR-05 & Dịch vụ AI không có đường công khai, yêu cầu khóa nội bộ & Không kết nối được cổng 8002 từ máy chủ; thiếu, sai, đúng khóa: 401, 401, 200 & Đạt \\
\hline
NFR-06 & HTTP chuyển 301 sang HTTPS; chỉ TLS 1.2, 1.3 & Mã 301; bắt tay TLS 1.2, 1.3 thành công; từ chối TLS 1.0, 1.1 chỉ kiểm tra qua cấu hình & Đạt \\
\hline
NFR-07 & Mọi API khai báo vai trò; sai vai trò 403 & Kiểm thử duyệt mọi route đạt; gọi chéo vai trò nhận 403; E2E UC2--UC5 không chạy hết & Đạt một phần \\
\hline
NFR-08 & Tài khoản bị khóa mất quyền ngay & Khóa: đăng nhập và token cũ 401; mở khóa: 200 & Đạt \\
\hline
NFR-09 & P95 kiểm tra tên $\leq 5$ s & P95 0,063 s (trực tiếp), 0,113 s (HTTPS); máy chủ thật: trung vị 0,136 s, lớn nhất 2,025 s & Đạt \\
\hline
NFR-10 & P95 kiểm tra một đề cương $\leq 30$ s & 20 tệp: P50 1,98 s, P95 4,94 s; báo cáo 350\,000 ký tự: 65,57 s & Đạt \\
\hline
NFR-11 & API nộp tệp $\leq 2$ s & P95 0,077 s, lớn nhất 0,091 s & Đạt \\
\hline
NFR-12 & P95 danh sách đề tài, tra cứu kho $\leq 1$ s & P95 0,018 s và 0,025 s; máy chủ thật: lớn nhất 0,266 s và 0,034 s & Đạt \\
\hline
NFR-13 & Đánh chỉ mục chỉ mã hóa bản ghi mới & Số lần gọi mô hình tăng đúng 1 & Đạt \\
\hline
NFR-14 & \texttt{utf8mb4\_\allowbreak vietnamese\_ci}, migration, chuẩn hóa NFC & 73 cột văn bản, 0 sai collation; NFC và NFD cho cùng kết quả & Đạt \\
\hline
NFR-15 & Client tĩnh, proxy \texttt{/api}, \texttt{ProxyTimeout} 120 s & Bản build Vite, \texttt{/api/health} qua proxy trả \texttt{ok} & Đạt \\
\hline
NFR-16 & Chrome và Edge trên máy tính & Chrome 1440 px: 35/35 màn hình; chưa kiểm thử tự động trên Edge & Đạt trên Chrome \\
\hline
NFR-17 & Ít nhất 100 SV, 30 GV & 580 SV, 70 GV & Đạt \\
\hline
NFR-18 & Ít nhất 200 đề tài đã đánh chỉ mục; 80 truy vấn có nhãn & 300 đề tài, cả 300 đã vào chỉ mục; 80 truy vấn & Đạt \\
\hline
NFR-19 & 50 tệp, 25 cặp, đủ 4 loại & 50 tệp, 25 cặp, 4 loại 7/7/6/5 & Đạt \\
\hline
\end{longtable}
}

Bộ nhớ của container dịch vụ AI đạt đỉnh 1\,364 MiB trong vòng đo NFR-10, dưới trần 2\,048 MiB. Trên máy chủ thật, với 20 mẫu, giá trị lớn nhất là chặn trên của P95 nên các chỉ tiêu vẫn đạt; nguyên nhân của một lời gọi kiểm tra tên mất 2,025 s chưa xác định được, và các phép đo có ghi dữ liệu (NFR-08, NFR-10, NFR-11) chưa được chạy trên máy chủ.

%=====================================================================
\section{Tài liệu hướng dẫn sử dụng}\label{sec:c4-huong-dan}
%=====================================================================

Người dùng truy cập hệ thống bằng Chrome hoặc Edge trên máy tính, đăng nhập bằng mã sinh viên, mã giảng viên hoặc email cùng mật khẩu do Bộ môn cấp; trang tổng quan và thông báo dẫn thẳng tới màn hình cần xử lý. Sinh viên nộp đề cương theo mẫu tên tệp hiển thị trên trang (PDF hoặc DOCX, tối đa 20~MB), theo dõi phê duyệt, nộp báo cáo tiến độ, báo cáo cuối và xem kết quả bảo vệ. Giảng viên giao đề tài ở \textit{Sinh viên của tôi}, duyệt và nộp đề cương lên Bộ môn ở \textit{Duyệt đề cương}, nhận xét tiến độ và chấm điểm. Bộ môn chuẩn bị đợt, quản lý tài khoản, phê duyệt đề cương, tổ chức bảo vệ, quản lý kho và ngưỡng. Các lệnh cài đặt và bảo trì cho quản trị viên được mô tả trong \path{deploy/README.md} của mã nguồn.

%=====================================================================
\section{Kết luận chương}
%=====================================================================

Chương này đã trình bày việc triển khai, kết quả thi công, kiểm thử và đánh giá hệ thống quản lý đồ án tốt nghiệp.

\textbf{Ưu điểm.}
\begin{itemize}
    \item Hệ thống bao phủ trọn quy trình đồ án của Bộ môn, từ phân công GVHD, giao đề tài, duyệt đề cương hai cấp, theo dõi tiến độ, tổ chức bảo vệ tới lưu kho, thay cho việc trao đổi qua email, tệp Excel và giấy tờ.
    \item Kiểm tra trùng lặp bằng AI chạy hoàn toàn trên máy chủ của Khoa. Mô hình được chọn và ngưỡng được hiệu chỉnh bằng số liệu, theo quy tắc viết trước khi đo, có tập kiểm chứng riêng; kết quả chỉ rõ đồ án nguồn và từng đoạn trùng để giảng viên tự đánh giá.
    \item Các chỉ tiêu thời gian đạt với khoảng dư lớn (kiểm tra đề cương P95 dưới 5 giây so với chỉ tiêu 30 giây); cài đặt bằng một lệnh, dưới 10 phút; bảo mật nhiều lớp từ HTTPS, phân quyền phía máy chủ tới tách dịch vụ AI và CSDL khỏi mạng ngoài.
\end{itemize}

\textbf{Hạn chế.}
\begin{itemize}
    \item \textit{Hệ thống không xác thực chữ ký trên đề cương}: khi phê duyệt, hệ thống chỉ kiểm tra bản tải lên là tệp PDF hợp lệ. Đây là lựa chọn có chủ đích để chấp nhận cả bản ký tay rồi quét lẫn bản ký bằng công cụ bên ngoài; trách nhiệm bảo đảm bản tải lên là bản đã ký thuộc về người phê duyệt.
    \item \textit{Ngưỡng được hiệu chỉnh trên dữ liệu tổng hợp}; danh sách phần mẫu chung là giả định. Bộ đánh giá kiểu PAN trên đồ án thật bù một phần, nhưng ca đạo văn ở đó vẫn là nhân tạo và hai bộ v1, v2 không phải phép so sánh có đối chứng. Ngưỡng cần được hiệu chỉnh lại khi có dữ liệu đồ án thật được gán nhãn.
    \item \textit{Phát hiện diễn đạt lại ở mức câu còn yếu}: chỉ 40--56\% câu diễn đạt lại được gắn cờ; đoạn do mô hình ngôn ngữ lớn viết lại đạt recall 0,563 ở bộ v1 và 0,793 ở bộ v2, chỉ nhờ tín hiệu ngữ nghĩa. Ngược lại, khi kho có nhiều đồ án cùng chủ đề, có thể có cảnh báo sai như cặp P19.
    \item \textit{Vận hành}: báo cáo dài kiểm tra lâu (khoảng 65 giây cho 350\,000 ký tự); người dùng chưa tự đổi được mật khẩu, chưa có sao lưu CSDL tự động, chỉ một máy chủ, chưa có giao diện di động.
    \item \textit{Kiểm thử chưa theo kịp thay đổi}: kịch bản E2E cho UC2--UC5 và kịch bản chụp màn hình chưa được cập nhật theo luồng duyệt hai cấp; trên máy chủ thật mới chạy các phép đo không ghi dữ liệu.
\end{itemize}

Các hạn chế trên là cơ sở cho hướng phát triển ở phần kết luận của khóa luận.

\end{document}
